\hypertarget{appendix-d52d5f78ffd1}{}
\section*{Carry and borrow}
\section{Carry and borrow}\label{carry-and-borrow}

\textbf{Curriculum depth:} Introduced · \textbf{First seen:} C002 · \textbf{Latest development:} C002

\subsection{Definition and mechanism}\label{definition-and-mechanism}

Carry is information leaving the high end of unsigned addition; borrow records an unmet subtraction at a bit position. Architecture subtraction flags may encode borrow or no-borrow.

For N-bit unsigned addition, carry-out means the exact sum exceeds 2ᴺ−1. A borrow means a subtraction needs value from a higher position. In an ordinary unsigned subtract, A\textless B implies borrow. Some architectures expose the complement, `no borrow', in their carry flag.

\subsection{Prerequisites}\label{prerequisites}

\begin{itemize}
\tightlist
\item
  \hyperref[appendix-83c3f28f7bda]{Positional notation} --- A base is the number of available digits.
\item
  \hyperref[appendix-24475e1596a8]{Boolean algebra} --- The logical domain is false/true.
\end{itemize}

\subsection{Used by}\label{used-by}

\begin{itemize}
\tightlist
\item
  \hyperref[appendix-d1a5b1e520c2]{Full adder} --- uses carry and borrow as part of its representation or reasoning.
\end{itemize}

\subsection{Recall and next depth}\label{recall-and-next-depth}

Explain carry and borrow without the label, work the example above, then state the assumption whose removal would change the result.

\textbf{Pending:} multi-precision arithmetic and architecture-specific conventions.

\subsection{Source and navigation}\label{source-and-navigation}

\href{Continuous\%20Notes/Computing\%20Foundations\%20-\%20From\%20Bits\%20to\%20C++\%20and\%20Systems\%20APIs\#Arithmetic\%20becomes\%20a\%20policy}{C002 teaching chapter} · \hyperref[appendix-fd090d982a15]{Part definitions and mechanisms} · \href{}{Dependency Map (02-Foundations-II, appendix-4ad47bdd269a)}
